\makerubrichead{Statement of Research}

\begingroup
\setlength{\parindent}{0pt}
\setlength{\parskip}{6pt}
\sloppy

My research programme connects \textbf{graph theory, networked systems, molecular and biomolecular graph modelling, and reliable scientific AI}. The common theme is structural reasoning: I study how mathematically defined graph structure can be converted into interpretable, testable, and reproducible evidence about complex systems.

A central strand of my work is \textbf{spectral and structural graph theory}, particularly graph energy, eccentricity-sensitive constructions, structural bounds, graph operations, and extremal behaviour. I developed and continue to study \textbf{Vertex Eccentricity Labeled Energy (VELE)}, an eccentricity-informed spectral graph invariant, together with related reconstruction, rigidity, and graph-construction questions. This theoretical work provides the mathematical foundation for several application-oriented directions.

A second strand concerns \textbf{network resilience and failure-sensitive analysis}. I investigate whether topology-derived quantities can provide useful pre-failure screening information in physical and infrastructure networks. Current work includes VELE-based screening of power-grid vulnerability validated against DC-flow behaviour, as well as leakage-controlled inspection prioritisation in flood-exposed road networks. A recurring principle in this work is to distinguish structural proxy value from physical or operational ground truth rather than treating graph metrics as automatically sufficient.

My interdisciplinary research also extends to \textbf{molecular and biomolecular graphs}. In EGFR QSAR, I study representation degeneracy and the generalization limits of classical graph descriptors under chemical-space shift. In protein stability, I investigate mutation-conditioned geodesic shielding in residue-contact networks, asking when alternative shortest-path structure can protect local network geometry from mutation. These projects combine graph-theoretic reasoning with reproducible computational experiments and explicit validity checks.

A fourth direction is \textbf{reliable scientific AI and evidence}. I work on applicability-gated scientific-rule learning, conformal prediction under composition shift, reliability under geographic transfer, explainable graph machine learning, and formal claim--evidence consistency. Across these projects, I emphasize leakage-controlled validation, uncertainty, calibration, external or grouped generalization, negative controls, transparent failure modes, and machine-checkable or reproducible evidence.

I maintain an open-research programme on GitHub so that code, frozen outputs, validation scripts, figures, and evidence boundaries remain inspectable. My long-term objective is to develop mathematically grounded methods that contribute both to graph theory and to reliable scientific computation across networks, chemistry, biology, infrastructure, and data-driven decision systems.

\endgroup
